\begin{figure}[htbp]
\centering
\includegraphics[width=0.75\textwidth]{figure_4}
\caption{
\textbf{Regional comparison of mean normalized uptake between real and MRI-synthetic PET images.}
(A,B) Violin plots show regional mean normalized uptake (arbitrary units) across cortical and subcortical regions for real and synthetic PET, comparing AV45 (A) and FDG (B). All regional uptake values were derived from per-image min--max-normalized PET intensities without anatomical reference-region normalization.
Panels A,B summarize the pooled available subjects without diagnostic stratification. Non-significant differences are not evidence of equivalence. (C--F) AD--CN regional comparisons. Violin plots display mean normalized uptake in AD (red) and CN (green) subjects for selected regions. Each region uses paired half-violins: left (solid) = real PET, right (dashed) = synthetic PET; colors match across modalities. Solid lines indicate real PET significance, dashed lines indicate synthetic PET significance. The statistical annotations in panels (C--F) follow the BH-FDR analysis: \texttt{****} corresponds to $q<10^{-4}$ and \texttt{ns} to $q\geq0.05$. (C,D) show representative regions with significant AD--CN differences for AV45 and FDG, whereas (E,F) show selected regions without significant AD--CN differences after FDR correction. Cblm-WM denotes Cerebellum-White-Matter. Additional regional agreement and effect-size analyses are provided in the Supplementary Materials.
}
\label{fig:4}
\end{figure}
